\section{Conclusion}
\label{sec:limit}

We studied on-policy distillation for long-horizon language agent tasks. Our
diagnosis shows that vanilla OPD suffers from two allocation mismatches:
full-horizon rollouts can spend compute on low-yield tail turns, while
trajectory-level KL reduction can concentrate loss on shallow tokens. TurnOPD
addresses these mismatches with turn-level budgeting: it adapts rollout depth
and progressively shifts KL normalization toward turn-balanced supervision.
Across ALFWorld, WebShop, and Multi-Hop Search, TurnOPD advances the
accuracy--time frontier over vanilla OPD.
Overall, our results suggest a broader lesson for agent OPD. The unit of
supervision in long-horizon agents should not be treated as a flat token
position, but as a turn-conditioned decision inside an evolving interaction
trace. TurnOPD opens a path toward turn-aware and budget-adaptive training of
long-horizon language agents.



% \section{Limitations}
% First, the evaluation covers three representative
% agent tasks, but it does not cover the full range of long-horizon workloads, such
% as GUI control, code agents, or open-ended deep research. Second,
% TurnOPD assumes access to a fixed teacher that can provide token-level distributions
% on student-visited states. This can be expensive when the teacher is large.
% Reducing teacher-query cost is a direction for future work.
